\section{Una línea, los pares y el resto de grado bajo}
Para una línea física $I$ de cualquier longitud $m\le j$, sus sitios son RC disjuntos.
Sean $S=\sum_I p_i\le M_0$, $Z_0=\Pr_{\rm Ber}(N_I\le1)\ge e^{-2M_0}$,
$x_i^0=(p_i/(1-p_i))/(1+\sum_I p_l/(1-p_l))$ y $d_i=x_i^0-p_i$.
La referencia de Bernoulli se usa para evaluar el polinomio de inclusión-exclusión.
Sus coeficientes para el cupo uno son $1,0,(-1)^{l-1}(l-1)$ para $l\ge2$.
Los dos truncados consecutivos en $L,L+1$ encajonan la probabilidad exacta por L03.
Con $L\sim A'\log j/\log\log j$, $2<A'<A$, su anchura es a lo sumo
$2L S^{L+1}/(L+1)!=O(S^2/j^2)$ por la parte anterior.
Si $m\le L$ no hay cola. Para $l\ge2$ la suma de productos es
$\le S^2 M_0^{l-2}/l!$. Así los errores relativos polinómicos, incluso tras
exponenciar, dan $O(S^2/j^2)$; el término lineal de grado uno está ausente.
En $D_i=p_i\partial_{p_i}$, $D_i Z_0/Z_0=d_i/(1-p_i)$, y para $i\ne h$,
\[
 \frac{D_iD_h Z_0}{Z_0}=\frac{d_i d_h-x_i^0x_h^0}{(1-p_i)(1-p_h)} .
\]
Insertando $\beta_i,G_{ih}=O(j^{-1})$ y $\sum|d_i|\le CS^2$, reemplazar esos
denominadores por uno cuesta $O(S^2/j^2)$. El cuadrado del primer cambio relativo
también tiene ese orden. Sale el término común de L02, líneas 24--32:
\[
 \log Z_0+\sum_i\beta_i d_i-\sum_{i<h}G_{ih}(d_i d_h-x_i^0x_h^0).
\]

Detallamos el cotejo canónico para identificar qué estimación importa. Las identidades
de L04--06 se transcriben completas con la prueba de sus restos. En el exterior
perfilado la acción residual $E$ cumple $E=E'=0$ en $p_I$ y
$E''_{ih}=2{\bf1}_{i=h}/(1-x_i)-G^{\rm ext}_{ih}$.
Retirar la masa del matching o mover las masas hasta $x^0$ cambia las entradas
cruzadas de $G$ en $O(j^{-2})$ contando los dos cruces mixtos. Las diagonales
de $E''$ cuestan $O(1/j)$ y $\sum d_i^2\le CS^3/j$.
El óptimo categórico $x^*$ está en $|x_i^*-x_i^0|\le Cp_iS^2/j$
por la contracción de las odds, de constante $CS/j$ (L04, 39--57).
La acción da $\log Z_0-\sum_{i<h}G_{ih}d_id_h+O(S^2/j^2)$.
En el prefactor RC el gradiente restringido a las marcas es
$(T_{II})^{-1}(T\nabla F_T)_I=\beta_I+O(j^{-2})$ por diagonal dominante
tras dividir cada fila por $p_i$. La Hessiana tangente tras restar las dos
incidencias comunes cuesta $O(j^{-1})$ en una RC y $O(j^{-2})$ en sitios disjuntos.
Las cargas de desplazamiento tienen total $O(S^2)$ y máximo $O(S/j)$ por fibra.
L05, líneas 13--48, da así el aporte $\sum\beta_i d_i+O(S^2/j^2)$.

La marginal $\beta_i$ entra en L03 y L05.
En el Schur de la línea, L06, líneas 3--28, la corrección matricial fuera de la
diagonal es $-t_i t_h/(1+\sum v_i)$, $t_i=\tau_3(x_i^*)$, $v_i=2(x_i^*)^2/(1-3x_i^*)$.
La parte diagonal y el logaritmo del factor escalar cuestan $O(S^2/j^2)$.
La matriz correctora tiene norma Hilbert--Schmidt $O(S/j)$ y
$\|G^y_{II}\|_{\rm op}=O(1)$ por ortogonalidad de las incidencias del matching.
El resto del logaritmo tras la traza es $O(S^2/j^2)$. Sustituir $t$ por $x^0$
y $G^y$ por $G$ da $+\sum_{i<h}G_{ih}x_i^0x_h^0$, con el mismo resto.
Queda $|e_I|\le CS^2/j^2$, también para $m=j$; para $m=0,1$ el pivote
unitario es exacto. En las dos familias $\sum_I S=2j$ y $\sum_I S^2\le2M_0 j$.

\paragraph{Una sola suma en el tablero.}
Para cuatro sitios en dos conflictos sobre líneas diferentes, las cotas individuales dan los pares
(dimensión geométrica, exponente del peso), en orden de cero a cuatro coincidencias RC:
\[
 (6,7),\quad(5,6),\quad(4,5),\quad(3,4),\quad(3,4).
\]
Las dos cuñas dan $(4,5),(3,4)$ según haya una RC común entre los extremos.
El primer par es M01. Los demás y las cuñas siguen la ruta alternativa certificada en el acta
W2\_r3, punto 4: los casos bosque, ciclo y varios ciclos de P05 y GI03--04, revisados junto con C02.
Los pesos T05, T07, T08 y T08c de la ronda 3 no se usan; sólo sirven de contraste.
La justificación siguiente se restringe a uniones sin otras relaciones diagonales.
En las restantes se usa la diferencia de cuatro casos del anexo T1.
Sólo se activan los conflictos elegidos y las restricciones RC; las incidencias
diagonales adicionales se mantienen apagadas al definir ese peso.
En las colocaciones con otras relaciones diagonales, la cota puntual de $w_{bd}$ es la de P05
para la activación restringida a $b$ y $d$, cuyo grafo es el de los dos conflictos más las
fibras RC: la prueba de P05 (GI02--04) sólo ve el exterior $A\setminus S$, las fibras RC
y los factores activados, y la diagonal no activada no entra en ningún subcaso.
Se necesitan los órdenes de las cotas, no los coeficientes exactos anteriores.
En el primer caso activar sólo los dos conflictos
crea un bosque en dos componentes si todas las RC son distintas, de donde $p=7$.
Esta segunda justificación no es independiente de M01: usa los mismos desarrollos
(P03, líneas 5--9), y la identidad de P04, líneas 54--63, reorganiza la cancelación
$+1/-1$ de MM06, líneas 26--36.
Con una coincidencia o más, el grafo de sitios y líneas, con filas y columnas, es conexo.
Una coincidencia une los dos árboles y da $p=6$, dos dan un ciclo y $p=5$,
y a partir de tres basta la cota genérica $p=4$ (P05, líneas 1--5, con las caras).
Una cuña tiene tres sitios y es árbol o tiene un ciclo; salen $p=5,4$.
Este argumento con una sola componente no exige tres conflictos ni conexión diagonal:
vale para dos conflictos unidos por fibras RC y para las cuñas (acta W2\_r3, líneas 214--221).
GI01, líneas 22--28, sólo enuncia soportes diagonalmente conexos, y P04 sólo hace falta
con dos o más componentes.
Las restricciones suprimidas se mantienen así en todas las caras de cada diferencia.
El coeficiente mixto, tomado sobre toda la unión, es su proyección de sitios:
al faltar un sitio desaparece uno de los conflictos y su diferencia es cero.
Para la dimensión se escogen coordenadas originales independientes (P06).
Dos pares en pendientes distintas tienen dos sitios base y dos desplazamientos.
Tres RC comunes fijan ambas filas de uno de los pares y una columna, o viceversa:
quedan a lo sumo tres parámetros. En la misma pendiente sobre líneas distintas,
el mismo recuento se aplica con el otro signo del desplazamiento; las coincidencias
que no pueden realizarse dan suma vacía. En la cuña una RC entre los extremos
impone igualdad de desplazamientos con signo. Cada tipo tiene a lo sumo $Cj^D$
colocaciones, aun con índices vivos no consecutivos. Por tanto su total es $O(1/j)$.

Para $v\le4$ el coeficiente del error completo en un soporte fijo tiene la cota
de P05, líneas 44--56, y P06 da $D\le p-1$. Se aplican primero a la proyección
sin sustracciones. Las líneas colineales ya se contaron enteras. En el resto sólo
se sustraen las parejas con unión exactamente ese soporte por inversión de Möbius;
dos componentes de dos sitios sin más conflictos ya quedan sustraídas.
Con tres o más conflictos hay 28 clases geométricas de cuatro sitios, contadas con sus
coincidencias RC (acta W2\_pares, líneas 148--165). En el camino y en la T,
cero, una y dos coincidencias dan $(D,p)=(5,6),(4,5),(3,4)$; en el 4-ciclo, cero y dos
coincidencias dependientes dan $(4,5),(3,4)$, y una sola no es realizable.
Se resta una pareja en el camino y en la T, y dos en el 4-ciclo; cada pareja restada
tiene exponente al menos $D+2$. En todas las clases $p=D+1$.
Hay un número absoluto de tipos. Se obtiene $O(1/j)$ para el resto de grado bajo.
Todas estas cotas se hacen en la máscara determinista de márgenes fijados.
Con $j=\lfloor bn\rfloor$ se pueden integrar directamente contra
$\mathbf1_{H_j}\,d\nu$. El cambio a $O_b(1/n)$ sólo ocurre en ese último paso.
